\section*{第 \ref{ch:b3:rovibrational-spectroscopy} 章——转动与振动光谱}

\begin{solution}{exo:b3:rovibrational-spectroscopy:1}
$I = h/(8\pi^2cB) = 6.62607 \times 10^{-34}/(8\pi^2 \times 2.99792 \times 10^{10} \times
1.9225) = \qty{1.4561e-46}{kg.m^2}$. $\mu = \qty{6.85621}{u} = \qty{1.13850e-26}{kg}$，故
$r_0 = \sqrt{I/\mu} = \qty{113.09}{pm}$。
\end{solution}

\begin{solution}{exo:b3:rovibrational-spectroscopy:2}
$2B_0(J+1)$：20.88、41.76、62.64、\qty{83.52}{cm^{-1}}。$J_{\max} = \sqrt{kT/2hcB} -
\frac12 = \sqrt{208.5/20.88} - 0.5 = 2.7$：能级 $J = 3$ 的
布居数最多，且最强谱线为 $4 \leftarrow 3$，位于 \qty{83.5}{cm^{-1}}。
\end{solution}

\begin{solution}{exo:b3:rovibrational-spectroscopy:3}
微波吸收需要永久偶极矩：仅 \ce{HCl} 和 \ce{H2O}。
转动拉曼需要各向异性极化率：\ce{H2}、\ce{HCl}、\ce{N2}、
\ce{CO2}、\ce{H2O}；球陀螺 \ce{CH4} 和 \ce{SF6} 除外。
\end{solution}

\begin{solution}{exo:b3:rovibrational-spectroscopy:4}
$N_1/N_0 = \eu^{-hc\tilde\nu/kT}$：在 \qty{300}{K} 时，$\eu^{-13.84} = \num{9.8e-7}$；在
\qty{1000}{K} 时，$\eu^{-4.15} = 0.016$。热带变得明显（约为基频的
百分之一），仅在 \qty{1000}{K} 附近。
\end{solution}

\begin{solution}{exo:b3:rovibrational-spectroscopy:5}
$\tilde\omega_e - 2\tilde\omega_ex_e = 2885.3$ 与 $2\tilde\omega_e - 6\tilde\omega_ex_e =
5665.0$：第二式减去第一式的两倍，得 $-2\tilde\omega_ex_e = -105.6$，
故 $\tilde\omega_ex_e = \qty{52.8}{cm^{-1}}$，且 $\tilde\omega_e = \qty{2990.9}{cm^{-1}}$。
\end{solution}

\begin{solution}{exo:b3:rovibrational-spectroscopy:6}
$D_e = 2990.95^2/(4 \times 52.82) = \qty{42341}{cm^{-1}}$；$G(0) = 1495.47 - 13.20 =
\qty{1482.3}{cm^{-1}}$；$D_0 = \qty{40859}{cm^{-1}} = \qty{488.8}{kJ/mol}$，与真实值
\qty{35760}{cm^{-1}} $= \qty{427.8}{kJ/mol}$ 相比，莫尔斯模型高估了 14\,\%。
\end{solution}

\begin{solution}{exo:b3:rovibrational-spectroscopy:7}
$R(0) - P(2) = 6B_0 = 62.64$：$B_0 = \qty{10.440}{cm^{-1}}$。$R(1) - P(1) = 6B_1 = 60.80$：
$B_1 = \qty{10.133}{cm^{-1}}$。$\alpha_e = \qty{0.307}{cm^{-1}}$。由 $R(0) = \tilde\nu_0 + 2B_1$ 得
$\tilde\nu_0 = \qty{2885.31}{cm^{-1}}$（与 $P(1) = \tilde\nu_0 - 2B_0$ 一致）。
\end{solution}

\begin{solution}{exo:b3:rovibrational-spectroscopy:8}
$\mu_H = \qty{0.97959}{u}$，$\mu_D = \qty{1.90441}{u}$，$\mu_H/\mu_D = 0.51438$。$B_0(\ce{DCl})
= 10.440 \times 0.51438 = \qty{5.370}{cm^{-1}}$。$\tilde\omega_e = 2990.95 \times \sqrt{0.51438}
= 2145.1$，$\tilde\omega_ex_e = 52.82 \times 0.51438 = 27.17$：$\tilde\nu_0 = 2145.1 - 54.3 =
\qty{2090.8}{cm^{-1}}$。
\end{solution}

\begin{solution}{exo:b3:rovibrational-spectroscopy:9}
位移为 $4B_0(J + \frac32)$：11.94、19.90、\qty{27.85}{cm^{-1}}（分别对应 $J = 0$、1、2）。两个
\ce{^{14}N} 核是自旋为 1 的全同玻色子：偶数 $J$ 能级
具有六种核自旋态，奇数 $J$ 能级为三种，因此偶数 $J$ 的谱线强度
为相邻谱线的两倍。
\end{solution}

\begin{solution}{exo:b3:rovibrational-spectroscopy:10}
刚性转子会使谱线位于 $\nu_1$、$2\nu_1$、$3\nu_1$：$230542.40$ 与
\qty{345813.61}{MHz}，分别比实测值高出 4.4 和 \qty{17.6}{MHz}。已知
$\nu_{J+1\leftarrow J} = 2B(J+1) - 4D(J+1)^3$：$\nu_1 = 2B - 4D$、$\nu_3 = 6B - 108D$，故
$3\nu_1 - \nu_3 = 96D$，得 $D = \qty{0.1835}{MHz}$，且 $B = (\nu_1 + 4D)/2 =
\qty{57635.97}{MHz}$。验算：$\nu_2 = 4B - 32D = \qty{230538.0}{MHz}$。
\end{solution}

\begin{solution}{exo:b3:rovibrational-spectroscopy:11}
$\Delta G(v + \frac12) = \tilde\omega_e - 2\tilde\omega_ex_e(v+1)$ 在 $v =
\tilde\omega_e/2\tilde\omega_ex_e - 1$ 附近为零；$n$ 项等差数列之和
从 $\tilde\omega_e - 2\tilde\omega_ex_e$ 递减至接近 0，其和约为 $n\tilde\omega_e/2 \approx
\tilde\omega_e^2/4\tilde\omega_ex_e - \tilde\omega_e/2 \approx D_e - G(0)$。对 \ce{HCl} 而言，28 个间隔
（$v = 0$ 至 27）之和为 \qty{40857}{cm^{-1}}，与 $D_e - G(0) = \qty{40859}{cm^{-1}}$ 相比。
\end{solution}

\begin{solution}{exo:b3:rovibrational-spectroscopy:12}
若只含偶数 $J$，跃迁 $J \to J + 2$ 的起点为 $J = 0, 2, 4, \dots$：位移为
$6B, 14B, 22B, \dots$，间距为 $8B$。\ce{CO2} 具有对称中心，且无
永久偶极矩，故无微波光谱。
\end{solution}

\begin{solution}{pb:b3:rovibrational-spectroscopy:1}
\textbf{1.} $B_0 = \nu/2 = \qty{57.6356}{GHz}$；$B_0/c = \qty{1.92252}{cm^{-1}}$。
\textbf{2.} $\mu = 12 \times 15.99491/27.99491 = \qty{6.85621}{u} = \qty{1.13850e-26}{kg}$。
\textbf{3.} $I = h/8\pi^2B_0 = \qty{1.4560e-46}{kg.m^2}$。
\textbf{4.} $r_0 = \sqrt{I/\mu} = \qty{113.09}{pm}$。
\textbf{5.} $r_0$ 比 $r_e = \qty{112.83}{pm}$ 大 \qty{0.26}{pm}：$B_0$ 是对 $1/r^2$ 在
零点振动上取平均的结果。在非谐势阱中，振动的核间距偏向
较大处，即在较长距离处停留的时间更多。
\textbf{6.} $2\nu_{1\leftarrow0} = \qty{230.5424}{GHz}$；实测谱线低 \qty{4.4}{MHz}：
离心畸变。
\textbf{7.} $0$、3.845 与 \qty{11.535}{cm^{-1}}，即 $E/k = 0$、5.53 与 \qty{16.60}{K}。
\textbf{8.} $N_2/N_1 = \frac53\exp[-(16.60 - 5.53)\,\mathrm K/T]$。
\textbf{9.} $\ln(\frac53/0.85) = 0.673 = 11.06/T$：$T = \qty{16}{K}$。
\textbf{10.} $J_{\max} = \sqrt{16.4/(2 \times 1.43878 \times 1.9225)} - 0.5 = 1.2$：$J = 1$。
\textbf{11.} 一个振动量子对应 $hc\tilde\nu/k \approx \qty{3080}{K}$：在
\qty{16}{K} 时，没有分子被振动激发，因此云团不发出振动
光。
\textbf{12.} 它只需 \qty{5.5}{K} 的激发能，即使在最冷的
气体中也能被激发，且 CO 丰度高，并具有偶极矩。
\textbf{13.} $\mu = 13.00335 \times 15.99491/28.99826 = \qty{7.17241}{u}$。
\textbf{14.} $B \propto 1/\mu$：$115.2712 \times 6.85621/7.17241 = \qty{110.1893}{GHz}$。
\textbf{15.} 预测值比实测谱线低 \qty{12}{MHz}（0.011\,\%）：$B_0 = B_e - \frac12\alpha_e$，且
振动–转动项 $\alpha_e$ 随 $\mu$ 的不同幂次而变化；
平衡结构相同，但零点平均并不相同。
\textbf{16.} $98.9/1.1 = 90$。
\textbf{17.} \ce{^{12}C^{16}O} 谱线已饱和（光学厚）：它不再
随气体的量增加而增强。
\textbf{18.} 稀有的 \ce{^{13}C^{16}O} 谱线仍与气体的量成正比，
而 \ce{^{12}C^{16}O} 谱线给出温度：二者合用，可同时测得二者。
\textbf{19.} $\tilde\nu_0 = 2169.81 - 2 \times 13.288 = \qty{2143.24}{cm^{-1}}$，\qty{4.666}{\micro m}。
\textbf{20.} $2\tilde\omega_e - 6\tilde\omega_ex_e = \qty{4259.89}{cm^{-1}}$。
\textbf{21.} $G(0) = \frac12\tilde\omega_e - \frac14\tilde\omega_ex_e = \qty{1081.58}{cm^{-1}}$。
\textbf{22.} $R(0) - P(1) = 2(B_0 + B_1) \approx 4B \approx \qty{7.7}{cm^{-1}}$。
\textbf{23.} \ce{^{13}C^{16}O}：约化质量更大，频率更低。
\textbf{24.} \textbf{$r_0(\ce{CO}) = \qty{113.1}{pm}$}，由单个射频数据得到。
\end{solution}
